\slajdid{04-s005}
\begin{frame}[label=04-s005]{Prelazak iz jednog u drugi brojni sistem}
\gusttekst
\alert{Prelazak iz brojnog sistema sa osnovom $r$ u decimalni brojni sistem}

Ako je prikazan broj
\[c_{n-1}c_{n-2}\,c_{n-3}\ldots c_1\,c_0.\,c_{-1}\,c_{-2}c_{-3}\ldots c_{-k+1}\,c_{-k}\]
u brojnom sistemu sa osnovom $r$, što najčešće zapisujemo
\[\left(c_{n-1}c_{n-2}\,c_{n-3}\ldots c_1\,c_0.\,c_{-1}\,c_{-2}c_{-3}\ldots c_{-k+1}\,c_{-k}\right)_r\]
Prelazak u decimalni brojni sistem (osnova 10 što je prirodno za nas) je direktno po izrazu
\[V=\sum_{i=-k}^{n-1}c_i r^i\]
Pri čemu se stepenovanje, množenje, sabiranje izvodi u decimalnom brojnom sistemu
\end{frame}
